% GENERATED FILE. Edit canonical lessons, then run npm run generate:books.
% canonical-source-hash: fnv1a64:6dfd1819de65e134
% canonical-lessons: FA-C421-buq, FA-C421-buq-zadan, FA-C421-cheragh, FA-C421-barfpakkon, FA-C421-ayene-baghal

\chapter{Horn, lights and mirrors}
\label{ch:fa-horn-lights-and-mirrors}
\hlchaptermodality{421}

\begin{chapteropening}
\textbf{By the end of this chapter:} \emph{I can talk about the horn, honking, the lights, the wipers and the side mirror.}

Complete the last lesson of chapter 421: I can talk about the horn, honking, the lights, the wipers and the side mirror.
\end{chapteropening}

\section[\texorpdfstring{\fa{بوق} (buq)}{بوق (buq)}]{\fa{بوق} (buq) — a horn (of a car)}

\label{lesson:FA-C421-buq}



\begin{quote}
Before the new one: say the Persian for to change gear, then the Persian for a clutch.
\end{quote}

\subsection*{You'll want to know: \fa{بوق}}
\textbf{\fa{بوق}} (\emph{buq}) — \textquotedblleft{}a horn (of a car)\textquotedblright{}.

\begin{quote}
\fa{بوق} \fa{این} \fa{ماشین} \fa{خیلی} \fa{بلند} \fa{است}.
\end{quote}

\begin{quote}
\emph{Buq-e in mâshin kheyli boland ast.}
\end{quote}

\begin{quote}
This car's horn is very loud.
\end{quote}

Also an old trumpet made of horn or brass.

\begin{cousinweb}
\emph{buq} was borrowed from Arabic and is usually traced to Arabic \emph{būq} 'a trumpet'. EtymDB does not list the word, so the line is hedged.
\end{cousinweb}

\begin{grammarlens}[title={What you've built}]
One more word for telling, explaining and describing at B1.
\end{grammarlens}

\subsection*{Your turn}
\begin{itemize}
\raggedright
  \item \emph{Say it:} \emph{buq}
  \item \emph{Say it:} \emph{buq}, once more
  \item \emph{Say it:} a sentence of your own with \emph{buq}
  \item \emph{Recall:} say the Persian for to change gear, then the Persian for a clutch, then say \emph{buq} again
\end{itemize}

\begin{tcolorbox}[breakable,colback=teal!4,colframe=teal!35!black,title={Before you move on}]
What does \fa{بوق} mean? (\textquotedblleft{}A horn (of a car)\textquotedblright{}.) Say it once more.
\end{tcolorbox}
\section[\texorpdfstring{\fa{بوق} \fa{زدن} (buq zadan)}{بوق زدن (buq zadan)}]{\fa{بوق} \fa{زدن} (buq zadan) — to honk, to sound the horn}

\label{lesson:FA-C421-buq-zadan}



\begin{quote}
Before the new one: say the Persian for a clutch, then the Persian for a horn.
\end{quote}

\subsection*{You'll want to know: \fa{بوق} \fa{زدن}}
\textbf{\fa{بوق} \fa{زدن}} (\emph{buq zadan}) — \textquotedblleft{}to honk, to sound the horn\textquotedblright{}.

\begin{quote}
\fa{راننده} \fa{تاکسی} \fa{جلوی} \fa{خانه} \fa{ما} \fa{دو} \fa{بار} \fa{بوق} \fa{زد}.
\end{quote}

\begin{quote}
\emph{Rânande-ye tâksi jolow-ye khâne-ye mâ do bâr buq zad.}
\end{quote}

\begin{quote}
The taxi driver honked twice in front of our house.
\end{quote}

\emph{buq} + \emph{zadan} (to strike). Present stem \emph{zan-}: \emph{buq mizanad} — he honks.

\begin{cousinweb}
\emph{buq zadan} is built from \emph{buq} 'horn', a loan usually traced to Arabic \emph{būq} 'a trumpet', and the native verb \emph{zadan} 'to strike'.
\end{cousinweb}

\begin{grammarlens}[title={What you've built}]
One more word for telling, explaining and describing at B1.
\end{grammarlens}

\subsection*{Your turn}
\begin{itemize}
\raggedright
  \item \emph{Say it:} \emph{buq zadan}
  \item \emph{Say it:} \emph{buq zadan}, once more
  \item \emph{Say it:} a sentence of your own with \emph{buq zadan}
  \item \emph{Recall:} say the Persian for a clutch, then the Persian for a horn, then say \emph{buq zadan} again
\end{itemize}

\begin{tcolorbox}[breakable,colback=teal!4,colframe=teal!35!black,title={Before you move on}]
What does \fa{بوق} \fa{زدن} mean? (\textquotedblleft{}To honk, to sound the horn\textquotedblright{}.) Say it once more.
\end{tcolorbox}
\section[\texorpdfstring{\fa{چراغ} (cherâgh)}{چراغ (cherâgh)}]{\fa{چراغ} (cherâgh) — a light, a lamp}

\label{lesson:FA-C421-cheragh}



\begin{quote}
Before the new one: say the Persian for a horn, then the Persian for to honk.
\end{quote}

\subsection*{You'll want to know: \fa{چراغ}}
\textbf{\fa{چراغ}} (\emph{cherâgh}) — \textquotedblleft{}a light, a lamp\textquotedblright{}.

\begin{quote}
\fa{شب} \fa{چراغهای} \fa{ماشین} \fa{را} \fa{روشن} \fa{میکنیم}.
\end{quote}

\begin{quote}
\emph{Shab cherâgh-hâ-ye mâshin râ roshan mikonim.}
\end{quote}

\begin{quote}
At night we switch on the car's lights.
\end{quote}

\emph{Cherâgh-e qermez} — a red light; \emph{cherâgh-e râhnamâ} — an indicator.

\begin{cousinweb}
\emph{cherâgh} is inherited: EtymDB traces it to Middle Persian \emph{čirāgh} 'lamp' and also links it with a Parthian form of the same word.
\end{cousinweb}

\begin{grammarlens}[title={What you've built}]
One more word for telling, explaining and describing at B1.
\end{grammarlens}

\subsection*{Your turn}
\begin{itemize}
\raggedright
  \item \emph{Say it:} \emph{cherâgh}
  \item \emph{Say it:} \emph{cherâgh}, once more
  \item \emph{Say it:} a sentence of your own with \emph{cherâgh}
  \item \emph{Recall:} say the Persian for a horn, then the Persian for to honk, then say \emph{cherâgh} again
\end{itemize}

\begin{tcolorbox}[breakable,colback=teal!4,colframe=teal!35!black,title={Before you move on}]
What does \fa{چراغ} mean? (\textquotedblleft{}A light, a lamp\textquotedblright{}.) Say it once more.
\end{tcolorbox}
\section[\texorpdfstring{\fa{برف} \fa{پاک} \fa{کن} (barfpâkkon)}{برف پاک کن (barfpâkkon)}]{\fa{برف} \fa{پاک} \fa{کن} (barfpâkkon) — a windscreen wiper}

\label{lesson:FA-C421-barfpakkon}



\begin{quote}
Before the new one: say the Persian for to honk, then the Persian for a light.
\end{quote}

\subsection*{You'll want to know: \fa{برف} \fa{پاک} \fa{کن}}
\textbf{\fa{برف} \fa{پاک} \fa{کن}} (\emph{barfpâkkon}) — \textquotedblleft{}a windscreen wiper\textquotedblright{}.

\begin{quote}
\fa{در} \fa{باران} \fa{برف} \fa{پاک} \fa{کن} \fa{تند} \fa{کار} \fa{میکند}.
\end{quote}

\begin{quote}
\emph{Dar bârân barfpâkkon tond kâr mikonad.}
\end{quote}

\begin{quote}
In the rain the wiper works fast.
\end{quote}

Literally 'snow-cleaner': \emph{barf} 'snow' + \emph{pâk} 'clean' + \emph{kon} 'doing'.

\begin{cousinweb}
\emph{barfpâkkon} is a native compound of \emph{barf} 'snow', inherited from Middle Persian \emph{wafr} (EtymDB), \emph{pâk} 'clean', from Middle Persian \emph{pāk} (EtymDB), and \emph{kon}, the present stem of \emph{kardan} 'to do'.
\end{cousinweb}

\begin{grammarlens}[title={What you've built}]
One more word for telling, explaining and describing at B1.
\end{grammarlens}

\subsection*{Your turn}
\begin{itemize}
\raggedright
  \item \emph{Say it:} \emph{barfpâkkon}
  \item \emph{Say it:} \emph{barfpâkkon}, once more
  \item \emph{Say it:} a sentence of your own with \emph{barfpâkkon}
  \item \emph{Recall:} say the Persian for to honk, then the Persian for a light, then say \emph{barfpâkkon} again
\end{itemize}

\begin{tcolorbox}[breakable,colback=teal!4,colframe=teal!35!black,title={Before you move on}]
What does \fa{برف} \fa{پاک} \fa{کن} mean? (\textquotedblleft{}A windscreen wiper\textquotedblright{}.) Say it once more.
\end{tcolorbox}
\section[\texorpdfstring{\fa{آینه} \fa{بغل} (âyene-baghal)}{آینه بغل (âyene-baghal)}]{\fa{آینه} \fa{بغل} (âyene-baghal) — a side mirror (of a car)}

\label{lesson:FA-C421-ayene-baghal}



\begin{quote}
Before the new one: say the Persian for a light, then the Persian for a windscreen wiper.
\end{quote}

\subsection*{You'll want to know: \fa{آینه} \fa{بغل}}
\textbf{\fa{آینه} \fa{بغل}} (\emph{âyene-baghal}) — \textquotedblleft{}a side mirror (of a car)\textquotedblright{}.

\begin{quote}
\fa{آینه} \fa{بغل} \fa{ماشین} \fa{من} \fa{شکست}.
\end{quote}

\begin{quote}
\emph{Âyene-baghal-e mâshin-e man shekast.}
\end{quote}

\begin{quote}
The side mirror of my car broke.
\end{quote}

\emph{âyene} 'mirror' + \emph{baghal} 'side'. The mirror inside the car is \emph{âyene-ye vasat}.

\begin{cousinweb}
\emph{âyene-baghal} is built from native \emph{âyene} 'mirror' and \emph{baghal} 'side, armpit', whose origin is not settled. EtymDB lists neither \emph{âyene} with older forms nor the Persian \emph{baghal}, so both parts are hedged.
\end{cousinweb}

\begin{grammarlens}[title={What you've built}]
One more word for telling, explaining and describing at B1.
\end{grammarlens}

\subsection*{Your turn}
\begin{itemize}
\raggedright
  \item \emph{Say it:} \emph{âyene-baghal}
  \item \emph{Say it:} \emph{âyene-baghal}, once more
  \item \emph{Say it:} a sentence of your own with \emph{âyene-baghal}
  \item \emph{Recall:} say the Persian for a light, then the Persian for a windscreen wiper, then say \emph{âyene-baghal} again
\end{itemize}

\begin{tcolorbox}[breakable,colback=teal!4,colframe=teal!35!black,title={Before you move on}]
What does \fa{آینه} \fa{بغل} mean? (\textquotedblleft{}A side mirror (of a car)\textquotedblright{}.) Say it once more.
\end{tcolorbox}
